% paper4.tex --- Arithmetic Landscape Theory, paper 4.
% Scaffolding (round 100). Section 2 is written; everything else is TODO-tagged.
% Every statement carries its status explicitly.
\documentclass[11pt,a4paper]{article}
\usepackage[T1]{fontenc}
\usepackage{lmodern}
\usepackage{amsmath,amssymb,amsthm}
\usepackage{mathtools}
\usepackage{booktabs}
\usepackage[margin=28mm]{geometry}
\usepackage{microtype}
\usepackage[colorlinks,citecolor=blue,linkcolor=blue,urlcolor=blue]{hyperref}
\usepackage{xcolor}

\theoremstyle{plain}
\newtheorem{theorem}{Theorem}[section]
\newtheorem{proposition}[theorem]{Proposition}
\newtheorem{lemma}[theorem]{Lemma}
\newtheorem{corollary}[theorem]{Corollary}
\theoremstyle{definition}
\newtheorem{definition}[theorem]{Definition}
\newtheorem{problem}[theorem]{Problem}
\newtheorem{conjecture}[theorem]{Conjecture}
\theoremstyle{remark}
\newtheorem{remark}[theorem]{Remark}

\newcommand{\lm}{\mathrm{lm}}
\newcommand{\Lean}[1]{\texttt{#1}}
\newcommand{\TODO}[1]{\textcolor{red}{[\textbf{TODO} #1]}}
\newcommand{\STATUS}[1]{\textcolor{blue}{\textsf{[#1]}}}

\title{Subset-sum landscapes under deformed and random measures:\\
the biased invariant $\Gamma^{(q)}$ and unconditional rigidity\\[4pt]
\normalsize\textsf{(working title --- scaffolding, round 100)}}
\author{Kentaro Amauchi\thanks{Independent researcher.}}
\date{\today}

\begin{document}
\maketitle

\begin{abstract}
\TODO{written last.} Papers 1--3 fixed both the sequence and the measure. This paper moves
the measure along two axes --- bias the coin, and randomise the sequence --- and collects the
two open questions recorded in paper~3.
\end{abstract}

\section{Introduction}\label{sec:intro}
\TODO{full introduction.} The two deformation axes, and why each is worth taking:
\begin{itemize}
\item \emph{Bias.} The classification theorem of paper~1 is a statement about which elements
      lie on which side of a target; it says nothing about how configurations are weighted.
      So the whole combinatorial layer survives verbatim when the uniform measure is replaced
      by the Bernoulli$(q)$ product measure, and the gap series deforms into an explicit
      one-parameter family. \S\ref{sec:biased} does this, and the answer has a shape worth
      stating in the introduction: \textbf{the fair coin gives the least rugged landscape}.
\item \emph{Randomness.} For a random odd sequence the minor-arc input that cost paper~2
      fifty rounds is replaced by a second-moment bound. \S\ref{sec:random} is therefore
      expected to carry \emph{no ineffective constant at all} --- the first such statement in
      the series. \TODO{write; see \S\ref{sec:random}.}
\end{itemize}

\section{The biased invariant}\label{sec:biased}

Let $P_q$ be the product measure under which each $a_i\in A$ is present independently with
probability $q\in(0,1)$, and write $\lm_q(n)$, $\deg_q(n)$ for the $P_q$-measures of the sets
of strict local minima and of ground states at the target $n$. The target is taken at the
$q$-tilted centre $n_q=\lfloor qT\rceil$.

\begin{definition}[biased gap series]\label{def:gq}
\[
  \Gamma^{(q)}(A)\;=\;1+\sum_{d\ge1}\Bigl[q^{N_d}+(1-q)^{N_d}\Bigr].
\]
\end{definition}

At $q=\tfrac12$ every bracket is $2\cdot2^{-N_d}$ and $\Gamma^{(1/2)}$ is the window series
$W_D$ of paper~1, hence $\Gamma(A)+(2D+1)2^{-k}$.

\begin{remark}[why the classification survives the deformation]\label{rem:measurefree}
Paper~1's classification is measure-free: at error $d$, an overshooting configuration is a
strict local minimum exactly when every $a\le2d$ is \emph{excluded}, and an undershooting one
exactly when every $a\le2d$ is \emph{included}. Under $P_q$ those two events have probability
$(1-q)^{N_d}$ and $q^{N_d}$, which is Definition~\ref{def:gq}. Writing
$r^{q}_B(m)=\sum_{U\subseteq B,\ \sigma(U)=m}q^{|U|}(1-q)^{|B|-|U|}$ and factorising $P_q$
across $I_d$ and $B_d$,
\[
  \mathrm{over}_q=(1-q)^{N_d}r^{q}_{B_d}(n+d),\quad
  \mathrm{under}_q=q^{N_d}r^{q}_{B_d}(n-d-\sigma_d),\quad
  \deg_q=\!\!\sum_{J\subseteq I_d}\!\!q^{|J|}(1-q)^{N_d-|J|}r^{q}_{B_d}(n-\sigma(J)),
\]
and the binomial sum over $J$ is exactly $1$, so a flat $r^{q}$ gives the $d$-th term of
Definition~\ref{def:gq}. \STATUS{derived; the assignment verified per stratum, see
Remark~\ref{rem:assign}}
\end{remark}

\subsection{The fair coin is the flattest}

\begin{lemma}[termwise minimality]\label{lem:termwise}
For every integer $N\ge0$ the function $f_N(q)=q^{N}+(1-q)^{N}$ on $(0,1)$ is symmetric about
$q=\tfrac12$ and convex, and $f_N(q)\ge f_N(\tfrac12)=2^{1-N}$, with strict inequality for
$q\ne\tfrac12$ as soon as $N\ge2$.
\end{lemma}

\begin{proof}
$f_N(1-q)=f_N(q)$ is immediate. For $N\ge2$,
$f_N'(q)=N\bigl[q^{N-1}-(1-q)^{N-1}\bigr]$ and
$f_N''(q)=N(N-1)\bigl[q^{N-2}+(1-q)^{N-2}\bigr]>0$, so $f_N$ is strictly convex; and
$t\mapsto t^{N-1}$ is strictly increasing, so $f_N'<0$ on $(0,\tfrac12)$ and $f_N'>0$ on
$(\tfrac12,1)$. For $N\in\{0,1\}$, $f_N$ is the constant $2$ or $1$.
\end{proof}

\begin{remark}[formally verified]\label{rem:lean22}
Lemma~\ref{lem:termwise} is machine-checked. \Lean{termwise\_min} gives the inequality for
every $N$ and every $q\in[0,1]$, \Lean{termwise\_min\_strict} its strictness for $N\ge2$ and
$q\ne\tfrac12$, and \Lean{termwise\_min\_iff} the equality case as a biconditional; the
normalisation $2\cdot(\tfrac12)^N=2^{1-N}$ is \Lean{two\_mul\_half\_pow}. The non-strict
bound is proved twice, once from convexity of $x^N$ on $[0,\infty)$ and once from Bernoulli's
inequality (\Lean{termwise\_min'}), so the two proofs check the \emph{statement} against each
other and not merely the tactic script. No \texttt{sorry}, and no axioms beyond
\texttt{propext}, \texttt{Classical.choice}, \texttt{Quot.sound}.
\STATUS{Lean-verified; \texttt{20260809\_TermwiseMin.lean}}
\end{remark}

\begin{theorem}[the fair coin minimises the invariant]\label{thm:fair}
For every finite increasing sequence $A$ of odd positive integers with $|A|\ge2$ and every
$q\in(0,1)$,
\[
  \Gamma^{(q)}(A)\;\ge\;\Gamma^{(1/2)}(A),
\]
with equality if and only if $q=\tfrac12$.
\end{theorem}

\begin{proof}
Sum Lemma~\ref{lem:termwise} over $d$. For strictness it suffices that some $d\le D$ has
$N_d\ge2$: take any $d\ge(a_2-1)/2$, which is at most $D=(a_k-1)/2$ because $a_2\le a_k$.
\end{proof}

\begin{remark}[what the theorem says]\label{rem:physics}
$q(1-q)$ is the per-element variance of the coin. Biasing the coin in either direction
lowers that variance and raises the invariant: \textbf{the more predictable the coin, the more
rugged the landscape}. The effect is exactly symmetric in $q\leftrightarrow1-q$, as it must
be, since complementing every configuration exchanges $q$ with $1-q$ and $n$ with $T-n$.
\end{remark}

\begin{corollary}[closed form for the odd numbers]\label{cor:oddsclosed}
For $A=(1,3,\dots,2k-1)$ one has $N_d=d$ for $d\le k$, so in the limit $k\to\infty$
\[
  \Gamma^{(q)}=1+\frac{q}{1-q}+\frac{1-q}{q}=\frac{1}{q(1-q)}-1\;\ge\;3,
\]
with equality only at $q=\tfrac12$: \emph{the invariant of the odd numbers is the reciprocal
of the coin's variance, minus one}. \STATUS{verified against the summed series to
$5.3\cdot10^{-15}$ over $q=0.10,\dots,0.80$; \texttt{e5\_r098}}
\end{corollary}

\subsection{Verification}\label{sec:e5}

\STATUS{experimentally confirmed} Measured $\lm_q/\deg_q$ at $n_q$ against the point
prediction of Definition~\ref{def:gq}, relative error:

\begin{center}
\begin{tabular}{llccc}
\toprule
profile & $q$ & $k=60$ & $k=100$ & $k=140$\\
\midrule
odds & 0.20 & $1.00\cdot10^{-2}$ & $3.16\cdot10^{-3}$ & $1.49\cdot10^{-3}$\\
odds & 0.30 & $2.63\cdot10^{-3}$ & $8.02\cdot10^{-4}$ & $3.75\cdot10^{-4}$\\
odds & 0.50 & $2.73\cdot10^{-4}$ & $5.99\cdot10^{-5}$ & $2.19\cdot10^{-5}$\\
primes & 0.30 & $2.18\cdot10^{-3}$ & $5.02\cdot10^{-4}$ & $2.20\cdot10^{-4}$\\
primes & 0.50 & $1.73\cdot10^{-4}$ & $2.26\cdot10^{-5}$ & $7.32\cdot10^{-6}$\\
\bottomrule
\end{tabular}
\end{center}

\noindent
There are no free parameters. Two features: the error decays faster than $1/k$ (fitted
exponent $2.25$ at $q=0.2$ to $3.0$ at $q=0.5$), and it grows monotonically as $q$ leaves
$\tfrac12$, by a factor $\approx40$ from $q=0.5$ to $q=0.2$ at fixed $k$. The second is what
the biased flatness hypothesis will have to quantify. \TODO{state the biased (H); design
after \S\ref{sec:e6}.}

\begin{remark}[the assignment, and a test that could not decide it]\label{rem:assign}
Definition~\ref{def:gq} is symmetric in $q\leftrightarrow1-q$, and so is the exact
complementation identity $\lm_q(n)=\lm_{1-q}(T-n)$ --- at the $q$-tilted centre,
$T-qT=(1-q)T$. Neither can therefore distinguish the assignment of $q^{N_d}$ and
$(1-q)^{N_d}$ to the two strata: measured at $k=100$, the ratio at $q=0.4$ and at $q=0.6$
agrees to every printed digit. The discriminating test is per stratum,
$\mathrm{over}_q/\deg_q$ against $(1-q)^{N_d}$ and $\mathrm{under}_q/\deg_q$ against
$q^{N_d}$; measured, both ratios lie in $[0.9978,1.0018]$ over $d\le6$, whereas a swap would
disagree by a factor of $161$ at $d=6$. \STATUS{experimentally confirmed; \texttt{e5\_r098}}
\end{remark}

\subsection{The bias breaks the reflection symmetry}\label{sec:e6}

\STATUS{experimentally confirmed, shape} Off the $q$-tilted centre the deviation
$\mathrm{dev}_q=(\lm_q/\deg_q)/\Gamma^{(q)}-1$ behaves differently at $q=\tfrac12$ and away
from it, and the reason is structural. Paper~3's transfer function is even in $\lambda$
because the two strata sit at $\pm\delta_d$ about a common centre --- the analytic form of
$\lm(n)=\lm(T-n)$. Under bias that symmetry maps $q$ to $1-q$ rather than to itself, so the
odd orders need not cancel, and the first-order coefficient is
\[
  L^{(q)}(A)\;=\;\sum_{d\ge1}\delta_d\Bigl[q^{N_d}-(1-q)^{N_d}\Bigr],
\]
which is antisymmetric in $q\leftrightarrow1-q$ and vanishes identically at $q=\tfrac12$.
Measured for the odd numbers at $k=80$: $L^{(q)}=-28.634,\,-8.843,\,0,\,+8.843,\,+28.634$ at
$q=0.3,0.4,0.5,0.6,0.7$, and for the primes $-83.717,\,-23.677,\,0,\,+23.677,\,+83.717$.

The data agree. Over the same range of targets, the spread (max over min) of
$|\mathrm{dev}_q-\mathrm{dev}_q(0)|$ divided by $\lambda_q$ and by $\lambda_q^2$:

\begin{center}
\begin{tabular}{llcc}
\toprule
profile & $q$ & divided by $\lambda_q$ & divided by $\lambda_q^{2}$\\
\midrule
odds & 0.30 & \textbf{1.329} & 9.644\\
odds & 0.50 & 5.091 & \textbf{1.001}\\
primes & 0.30 & \textbf{1.284} & 8.858\\
\bottomrule
\end{tabular}
\end{center}

\noindent
So the response is linear in $\lambda$ under bias and quadratic at $q=\tfrac12$, where
$(\mathrm{dev}-\mathrm{dev}_0)/\lambda^{2}$ is constant to four digits ($19.8186$ to
$19.8352$ across a fivefold range of the target offset). \STATUS{experimentally confirmed;
\texttt{e6\_r100}} \TODO{the primes at $q=\tfrac12$ are not a clean test in this range:
$\mathrm{dev}-\mathrm{dev}_0$ crosses zero inside it, and dropping the single point where it
is $2.7\cdot10^{-7}$ leaves the quadratic ratio in $[45.06,48.86]$. Re-measure on a target
range that does not straddle the crossing.}

\section{The biased transfer function}\label{sec:phiq}
\TODO{fable-5: derive $\Phi^{(q)}$ by the partition-function method of paper 3's
\textup{rem:mushift}, with the biased normaliser. The two strata carry unequal weights
$q^{N_d}$ and $(1-q)^{N_d}$ and a $q$-shifted $\mu_d$, so the $\cosh$ becomes a general
combination of two exponentials; \S\ref{sec:e6} says the surviving odd part has coefficient
$L^{(q)}$, which is the first thing the derivation must reproduce.}

\section{Random sequences, unconditionally}\label{sec:random}

Before the theorem, the floor it has to stand on. Let $A$ be drawn from the Cram\'er model on
the odd integers: each odd $m\in[3,N]$ is included independently with probability $2/\log m$,
and $N$ is fixed by $\mathbb{E}|A|=k$. Nothing is conditioned, so $|A|$ fluctuates from
instance to instance. That is deliberate: $kR_A(k)$ is the $k$-invariant quantity of paper~3's
\textup{rem:onek}, so an ensemble of genuinely different sizes tests the $1/k$ law for free.
Twenty instances at each of two sizes, each through the whole paper-3 pipeline with its own
$N_d$ and $\sigma_d$:

\begin{center}
\begin{tabular}{lccccc}
\toprule
$\mathbb{E}|A|$ & range of $k$ & (H) series & mean $kR$ & sd & sd/mean\\
\midrule
$120$ & $101$--$135$ & $133.79$--$320.79$ & $214.96$ & $10.58$ & $4.92\%$\\
$180$ & $146$--$191$ & $132.94$--$471.77$ & $206.32$ & $6.94$  & $3.37\%$\\
\bottomrule
\end{tabular}
\end{center}

\noindent
at $x=0.06$; $c_A=kR/100$. Three things, and the third is the one to keep.
\STATUS{experimentally confirmed; \texttt{e7\_r102}}

\emph{(i) The hypothesis holds on the ensemble.} All forty instances have a bounded (H)
series and a window that tracks the primes'. This is a precondition, not a result: a theorem
about random sequences is unconditional only if the random sequences satisfy (H) themselves.

\emph{(ii) The constant concentrates.} The spread across instances is under $5\%$ and falls
with $k$, while the mean of $kR$ moves by $4\%$ between the two sizes --- the $1/k$ law of
paper~3 surviving an ensemble it was not fitted on.

\emph{(iii) The primes sit inside the ensemble and the odd numbers do not.} Against the
$k\approx180$ ensemble ($2.0632\pm0.0694$ in $c_A$), the primes score $z=-0.39$ and the odd
numbers $z=-2.70$, at every one of the four targets. In this statistic the primes are
indistinguishable from a random odd sequence of the same density; the odd numbers are a
different object.

There is one place where the random model is \emph{not} like the primes, and it is the place
that matters for the minor arcs. Write
$h_d(q)=\max_{(r,q)=1}\bigl(\prod_{a\in B_d}|\cos(\pi ra/q)|\bigr)^{1/|B_d|}$.
Since $|\cos(\pi a/4)|=1/\sqrt2$ for \emph{every} odd $a$, the modulus-$4$ floor
$h_d(4)=1/\sqrt2$ is deterministic, not merely universal. Paper~2's vanishing lemma kills the
modulus-$2m$ peak ($m$ odd) as soon as $B_d$ contains an element $\equiv3\pmod 6$. A set of
primes contains exactly one such element, namely $3$, and $3$ leaves $B_d$ at $d\ge2$, so the
primes keep $h_d(6)=\sqrt3/2$; a random odd sequence contains about $|B_d|/3$ of them and
loses the peak at every $d$. Measured over $q\le120$ at $d=1,2,5$: every random instance has
$\max_{q\ge3}h_d(q)=1/\sqrt2$ attained at $q=4$, and the primes have $\sqrt3/2$ at $q=6$.
\STATUS{experimentally confirmed; \texttt{e7\_r102}}

\noindent
So the random theorem should come out with the geometric rate $1/\sqrt2=0.7071$ per element
rather than the $\sqrt3/2=0.8660$ of papers 1--2: \emph{the primes are the harder case, and
the modulus-$6$ ripple is a structural feature of the primes that the random model destroys.}

\TODO{fable-5: for random odd $A$ the minor-arc input is a second-moment bound
$\mathbb{E}|S_A|^{2}=b$ plus Borel--Cantelli on a $\theta$-net --- no Kumchev, no
Siegel--Walfisz, no Siegel constant. If it lands it is the trilogy's first
zero-ineffective-constant rigidity theorem. The measurements above are the floor it must
clear, and they sharpen the brief: the rate to aim for is $1/\sqrt2$, not $\sqrt3/2$.}

\section{Two questions inherited from paper 3}\label{sec:inherited}
\emph{(i) The $c_d$ fold-in is settled, and it goes paper 3's way.} The share of the
residual that $c_d$ carries falls like $k^{-2\alpha}$ --- measured exponents $1.98$, $2.35$,
$2.99$ against predicted $2.00$, ${\approx}2.4$, $3.00$ on the odd numbers, the primes and
$\alpha=\tfrac32$, with the squares already below the numerical floor. So $\Phi$ as printed
is right for every larger $k$ and $c_d$ belongs to $E_k$ permanently; paper 3's
\textup{rem:cd} now says so. \STATUS{experimentally confirmed, four profiles, $k\le300$;
\texttt{e8\_r103}}

\emph{(ii) $c_A=c_A'$ is still open, but sharper.} Carried to $k=300$ the ratio reaches
$1.0064$ (odd numbers), $1.0001$ (primes), $0.9984$ ($\alpha=\tfrac32$), and the gap closes
like $k^{-2.4}$ --- faster than the $1/k$ law both constants sit on.
\STATUS{experimentally confirmed; \texttt{e8\_r103}}
\TODO{the mechanism is still missing, and a measured rate of convergence is evidence about a
limit, not a proof of one. F45 stands: substituting $s$ for $\lambda$ makes the residual
three times worse.}

\section{Outlook}\label{sec:outlook}
\TODO{the algorithmic reading --- local search on structured instances --- which is where
this programme started. Outlook only; no claims.}

\section{Honest scope}\label{sec:scope}
\begin{itemize}
\item \emph{Lean-verified}: Lemma~\ref{lem:termwise} in all three forms
      (Remark~\ref{rem:lean22}), including the equality case.
\item \emph{Proved}: Lemma~\ref{lem:termwise}, Theorem~\ref{thm:fair},
      Corollary~\ref{cor:oddsclosed}; and Remark~\ref{rem:measurefree}'s derivation of
      Definition~\ref{def:gq} modulo the biased flatness hypothesis, which is not yet stated.
\item \emph{Experimentally confirmed, with the range stated}: the point prediction
      (\S\ref{sec:e5}, two profiles, five $q$, three $k$), the stratum assignment
      (Remark~\ref{rem:assign}), the linear response under bias (\S\ref{sec:e6}).
\item \emph{Not yet written}: \S\S\ref{sec:phiq}--\ref{sec:outlook}, and the biased
      hypothesis that \S\ref{sec:biased} needs to become a theorem.
\end{itemize}

\begin{thebibliography}{9}
\bibitem{P1} K.~Amauchi, \emph{The gap series of an integer sequence}. Manuscript, 2026.
\bibitem{P2} K.~Amauchi, \emph{Asymptotic flatness of subset-sum landscapes of primes}.
Manuscript, 2026.
\bibitem{P3} K.~Amauchi, \emph{The transfer function of subset-sum landscapes}.
Manuscript, 2026.
\end{thebibliography}

\end{document}
